% Target language: Modern Standard Hindi (India), hi-IN.
% Source: content/en-US.tex, checked against content/zh-CN.tex at upstream commit 0d064f6.
% Translation: AI-assisted; needs native-speaker review before publication.
\StoryParagraph{3}{एक अंडाणु का वज़न: 3.5 माइक्रोग्राम।}

\StoryParagraph{3}{नकद पाँच करोड़ डॉलर का वज़न: 2.5 टन।}

\StoryParagraph{3}{\textenglish{Montage Laguna Beach} से फ़ोन पर जिंग तियान ने पहली चीज़ को गिरवी रखकर मुझसे दूसरी माँगी।}

\StoryParagraph{3}{वह अपने अंडाणु निकलवाने वहाँ गई थी। सरोगेट माँ की कोख में पलने वाला बच्चा 2027 में आने वाला था। अगर सब कुछ काफ़ी तेज़ी से होता, तो वह अब भी घोड़े के वर्ष में जन्म ले सकता था। यह बात उसी ने छेड़ी थी। क्लिनिक में उसने कहा था: पाँच करोड़ नहीं, तो अंडाणु भी नहीं निकलेंगे।}

\UnprintedPause


\StoryParagraph{0}{\textenglish{Montage Laguna Beach} प्रशांत महासागर के ऊपर एक समुद्री चट्टान पर बना है। उसकी निजता के लिए हमने पूरी मंज़िल तीस दिनों के लिए, दस लाख डॉलर में, बुक कर ली। वह आठ दिन बाद चली गई।}

\StoryParagraph{0}{बाकी बाईस दिनों तक होटल के सफ़ाईकर्मी रोज़ पहले की तरह ऊपर जाते रहे। मेरी सहायक ने मुझे बताया। चाबी लौटाने गई तो उसने एक बार यह दृश्य देखा था: एक मैक्सिकन महिला ठेला धकेलते हुए कमरे-दर-कमरे जा रही थी, तौलिये बदल रही थी, काउंटर पोंछ रही थी, चादरें ठीक कर रही थी, बाहर निकलकर हर दरवाज़ा अपने पीछे बंद कर रही थी। पूरी मंज़िल साफ़ करने में दो घंटे लगते थे।}

\StoryParagraph{0}{मेरी सहायक ने उससे पूछा कि जब वहाँ कोई ठहरा ही नहीं है, तो वह अब भी यह सब क्यों करती है।}

\StoryParagraph{0}{उसने कहा, रजिस्टर में तो कोई ठहरा हुआ था।}

\StoryParagraph{0}{अपनी रानी के बिना वर्साय।}

\StoryParagraph{0}{तभी मुझे लगा कि जिंग तियान को यह पसंद था। उसे पसंद था कि उसकी वजह से कोई जगह खाली पड़ी रहे। कुछ भी न हो।}

\StoryParagraph{3}{\textenglish{Montage Laguna Beach} के साथ एक \textenglish{Gulfstream G700} भी खाली पड़ा था, \textenglish{Van Nuys} में अटपटा-सा खड़ा, जबकि उसका चालक-दल रोज़ पहले की तरह काम पर आता रहा।}

\StoryParagraph{3}{आठवें दिन उसने वह फ़ोन किया। उसके बाद किसी ने उसे फिर नहीं देखा।}

\NarrativeBreak


\StoryParagraph{0}{2007 तक कोई उसके लिए पैसे जला रहा था। उसी साल मैंने पेकिंग विश्वविद्यालय में पढ़ाई शुरू की थी।}

\StoryParagraph{3}{बेशक, पैसे जलाने वाला मैं नहीं था।}

\StoryParagraph{3}{\textenglish{QQ} ने एक पॉप-अप विंडो खोली, जिसमें वह विज्ञापन का चेहरा थी। वह मेरे पुराने \textenglish{Lenovo IBM} कंप्यूटर पर था; रिज़ॉल्यूशन इतना धुँधला था कि कुछ साफ़ नहीं दिखता था। मैंने तस्वीर सहेज ली।}

\StoryParagraph{3}{पेकिंग विश्वविद्यालय के अपने छात्रावास से, जहाँ एक साल का किराया एक हज़ार युआन था, मैं ऑनलाइन गया, उसे \textenglish{Xiaonei} पर खोजा और मित्रता अनुरोध भेजा। \textenglish{Xiaonei} में साथ एक संदेश देना ज़रूरी था। करीब चालीस मिनट तक मैंने लिखा और मिटाया, मिटाकर फिर लिखा। अंत में भेजा: “नमस्ते।”}

\StoryParagraph{3}{उसने अनुरोध स्वीकार नहीं किया।}

\StoryParagraph{3}{शायद उसने कभी \textenglish{Xiaonei} इस्तेमाल ही नहीं किया था।}

\StoryParagraph{3}{बाद के वर्षों में मैं जो भी करता, एक सवाल बार-बार मन में आता: मुझे कहाँ तक पहुँचना होगा कि वह स्वीकार करे? अपना पहला अरब कमाया, तब भी यह सोचा। दस अरब पर फिर सोचा।}

\StoryParagraph{3}{आख़िर मैंने इस बारे में सोचना छोड़ दिया। \textenglish{Renren} वर्षों पहले बंद हो चुका था।}

\StoryParagraph{3}{मैंने दूसरा अनुरोध कभी नहीं भेजा।}

\StoryParagraph{3}{2011 में बीजिंग में कड़ाके की ठंड थी। मैंने 150 युआन की डाउन जैकेट खरीदी। कीमत इसलिए याद है कि तीन दिन तक सोचता रहा था।}

\StoryParagraph{3}{उसी साल उसने पंद्रह करोड़ युआन की फ़िल्म \textenglish{\textit{The Warring States}} में मुख्य भूमिका निभाई। \textenglish{Sun Honglei}, \textenglish{Francis Ng}, \textenglish{Kim Hee-sun}, \textenglish{Kiichi Nakai} और \textenglish{Jiang Wu} को बस पृष्ठभूमि बनाने के लिए बुलाया गया था, ठीक वैसे ही जैसे उन्नीस साल बाद \textenglish{Montage Laguna Beach} और \textenglish{Gulfstream G700} को। उन्हें पता नहीं था कि वे वहाँ क्यों हैं। इतनी दूर आए थे, तो कुछ न कुछ होना चाहिए था। कुछ नहीं हुआ।}

\StoryParagraph{3}{फिर पंद्रह करोड़ डॉलर के बजट वाली \textenglish{\textit{The Great Wall}} आई। \textenglish{Zhang Yimou}। \textenglish{Andy Lau}। \textenglish{Matt Damon}। फिर \textenglish{Jackie Chan}। जिंग तियान के संसार की सूची बढ़ती चली गई।}

\StoryParagraph{3}{तब तक मैं संयुक्त राज्य अमेरिका में था, पढ़ रहा था, कंपनी शुरू कर रहा था और दिन-ब-दिन व्यस्त होता जा रहा था। केवल एक बात कभी नहीं बदली: जब भी डिवाइस बदलता, चाहे जितनी परेशानी हो, \textenglish{QQ} वाली वह तस्वीर साथ ले जाता।}

\StoryParagraph{3}{उन्नीस साल बाद मेरा नाम भी उस सूची में जुड़ गया।}

\NarrativeBreak

\StoryParagraph{3}{हम पहली बार \textenglish{Rosewood Hong Kong} में मिले। मैं हमेशा सोचता था कि जिस दिन उससे मिलूँगा, घबरा जाऊँगा।}

\StoryParagraph{3}{मैं नहीं घबराया।}

\StoryParagraph{3}{वह तस्वीरों से अधिक दुबली थी, सहज कपड़े पहने सोफ़े पर पीछे टिककर बैठी थी। मैंने उसे देखा और अचानक याद नहीं कर पाया कि इतने समय से ठीक किस चीज़ की प्रतीक्षा कर रहा था।}

\StoryParagraph{3}{मैं उसे \textenglish{Cattelan} का केला दिखाने ले गया। वह देर तक उसे देखती रही, फिर मुझसे पूछा, “यह कितने का है?”}

\StoryParagraph{3}{बासठ लाख डॉलर।}

\StoryParagraph{3}{उसने कहा, “जब यह सड़ जाएगा, तब क्या होगा?”}

\StoryParagraph{3}{मैंने कुछ नहीं कहा। उसने कहा, “छोड़ो। चलो फ़िल्म देखते हैं।”}

\StoryParagraph{3}{मैंने कहा, “ठीक है।”}

\StoryParagraph{3}{उसने कहा कि हॉल में हम दोनों के अलावा कोई नहीं होना चाहिए।}

\StoryParagraph{3}{मैंने कहा, “ठीक है।”}

\StoryParagraph{3}{मैं फ़ोन करने बाहर गया। लौटा तो वह अब भी केले के सामने खड़ी थी।}

\StoryParagraph{3}{फ़िल्म \textenglish{\textit{Zootopia 2}} थी।}

\StoryParagraph{3}{दो लोगों के लिए पूरा सिनेमा हॉल। ऐसा इंतज़ाम आम तौर पर तभी होता है जब केंद्रीय प्रचार विभाग किसी फ़िल्म की समीक्षा करता है।}

\StoryParagraph{3}{मैंने अपनी सहायक से सिनेमाघर की हर सीट खरीदवा ली।}

\StoryParagraph{3}{मैंने नकली बॉक्स ऑफिस पर प्रहार करने वाले लेख पढ़े थे: सारे टिकट बिके, हॉल में केवल दो लोग। अब समझ आया। शायद बात हमेशा बॉक्स ऑफिस की नहीं होती।}

\StoryParagraph{3}{सारी सीटें खरीदना कठिन नहीं था। समस्या वे टिकट थे जो पहले ही बिक चुके थे। मेरे सहायक और अंगरक्षक जल्दी पहुँचकर प्रवेश-द्वार पर इंतज़ार करने लगे। टिकट लेकर जो भी आता, वे उसे नकद पैसे दे देते।}

\StoryParagraph{3}{एक दंपती अपने बच्चे के साथ आया। बच्चा दरवाज़े तक पहुँच चुका था और पॉपकॉर्न का बड़ा डिब्बा हवा में ऊँचा उठाए था।}

\StoryParagraph{3}{एक सहायक ने नोटों की गड्डी आगे कर दी। बच्चा समझ नहीं पाया कि वे भीतर क्यों नहीं जा सकते। माता-पिता ने पैसे लिए, धन्यवाद कहा, कुछ नहीं पूछा और चले गए।}

\StoryParagraph{3}{उस दिन छब्बीस लोगों ने पैसे लिए।}

\StoryParagraph{3}{फ़िल्म के बाद उसने \textenglish{WeChat} की प्रोफ़ाइल तस्वीर बदलकर \textenglish{Officer Judy Hopps} की तस्वीर लगा ली।}

\StoryParagraph{3}{उस रात मैंने उसके स्तन सहलाए। उसने मुझे परे धकेल दिया।}

\StoryParagraph{3}{मुझे लगा कि वह ऐसा नहीं चाहती। मैंने हाथ पीछे खींच लिया और कुछ कहने ही वाला था।}

\StoryParagraph{3}{उसने कुछ नहीं कहा। उठकर बैठी, बैग से एक छोटी थैली निकाली और उसकी ज़िप खोली। भीतर चीज़ें ऐसे सजी थीं, जैसे कोई छोटी-सी औज़ार-पेटी हो।}

\StoryParagraph{3}{उसने कहा, “तुम्हारे नाखून चुभते हैं।”}

\StoryParagraph{3}{मैंने देखा। मैं हमेशा नेल कटर से जैसे-तैसे नाखून काट लेता था। किनारे नुकीले रह जाते थे। आम तौर पर मेरा ध्यान ही नहीं जाता था।}

\StoryParagraph{3}{उसने कहा, “इधर आओ।”}

\StoryParagraph{3}{उसने मेरा हाथ अपनी गोद में रखा और एक-एक नाखून घिसने लगी।}

\StoryParagraph{3}{पहले खुरदरी सतह। उससे काम पूरा हुआ तो उसने सतह बदली, फिर एक बार और बदली। बीच-बीच में वह उँगली को रोशनी में उठाकर देखती और फिर घिसने लगती।}

\StoryParagraph{3}{मैंने पूछा, “इसमें कितना समय लगेगा?”}

\StoryParagraph{3}{उसने कहा, “जब तक ये लोगों को खरोंचना बंद न कर दें।”}

\StoryParagraph{3}{दस उँगलियाँ। उसे लगभग एक घंटा लगा।}

\StoryParagraph{3}{काम करते हुए वह मुझे समझाती गई: इस सतह से आकार बनता है, इससे चमक आती है, और आख़िरी सतह को आगे-पीछे मत रगड़ना। केवल एक दिशा में चलाना।}

\StoryParagraph{3}{मैंने कहा, “ठीक है।”}

\StoryParagraph{3}{काम पूरा होने पर उसने मेरा हाथ पलटा, उसे देखा और अपनी उँगलियाँ उस पर फेरीं।}

\StoryParagraph{3}{उसने कहा, “अब ठीक है।”}

\StoryParagraph{3}{“अब जहाँ चाहो, छू लो।”}

\StoryParagraph{3}{एक समय उसने कुछ कहा। उसने कहा, “किसी दिन मैं चली गई, तो फिर कोई तुम्हारे नाखून चमकाकर नहीं देगा।”}

\StoryParagraph{3}{मैं मुस्कराया। कुछ नहीं कहा।}

\StoryParagraph{0}{उससे मिलने के तीन महीने बाद मैंने किसी से एक फ़ाइल तैयार करवाई। वर्षों के उसके प्रोजेक्ट, कारोबारी साझेदार, हिस्सेदारियाँ, \textenglish{Jing Guoqing}, \textenglish{Tian Xin'ai}, उनसे जुड़े सारे खाते, बीजिंग में किराए की जगह, शंघाई में खरीदी जगह और कुछ दूसरी बातें। एक सप्ताह लगा। कुल मिलाकर चालीस से कुछ अधिक पन्ने। उसमें मर्दों और औरतों के बारे में भी बातें थीं। मुझे लगा वे सबसे कम मायने रखती थीं।}

\StoryParagraph{0}{उसमें दर्ज कुछ बातें उसने बाद में मुझे बताईं। उसने जो बताया, वह वहाँ लिखी बात से पूरी तरह मेल नहीं खाता था। मैंने कुछ नहीं कहा।}

\StoryParagraph{0}{कुछ बातें उसने कभी नहीं बताईं। मैंने कभी नहीं पूछा।}

\StoryParagraph{0}{वह जो भी कहती, मैं उस पर विश्वास करता था।}

\StoryParagraph{0}{मैंने वह फ़ाइल संभालकर रखी। वह अब भी मेरे पास है।}

\StoryParagraph{3}{लेकिन मैंने उसके एक भी शब्द पर विश्वास नहीं किया।}

\NarrativeBreak

\StoryParagraph{3}{उसने कहा, “हांगकांग में बहुत भीड़ है। लोग मुझे पहचान लेंगे। कहीं और चलते हैं।”}

\StoryParagraph{3}{मैंने कहा, “ठीक है।”}

\StoryParagraph{3}{उसने नक्शे पर \textenglish{Tenerife} की ओर इशारा किया।}

\StoryParagraph{3}{मैंने कहा, “ठीक है।”}

\StoryParagraph{3}{मैं फ़ोन करने बाहर गया। लौटकर ही पता चला कि वह जगह अटलांटिक महासागर में, अफ़्रीका के उत्तर-पश्चिमी तट के पास थी। अगले दिन दोपहर में हमारा तीस लोगों का दल रवाना हुआ। हवाई अड्डे पर एक \textenglish{A330} प्रतीक्षा कर रहा था।}

\StoryParagraph{3}{विमान के शयनकक्ष में उसने मुझसे पूछा, “मेरा नाम सुनने में अच्छा लगता है?”}

\StoryParagraph{3}{मैंने कहा, “हाँ।”}

\StoryParagraph{3}{उसने कहा, “यह मेरे पिता और माँ के उपनामों को जोड़कर बना है। मज़ेदार संयोग है न?”}

\StoryParagraph{3}{मैंने कहा, “मुझे यह पहले ही दिन याद हो गया था।”}

\StoryParagraph{3}{उसने बाँहें मेरे गले में डाल दीं और कहा, “अब मुझे जिंग तियान मत कहना। मुझे मम्मी कहो।”}

\StoryParagraph{3}{“ठीक है, मम्मी।”}

\StoryParagraph{3}{उतरने के बाद कप्तान मेरे पास आया और बोला, “बॉस, पहले द्वीप हमें उतरने नहीं दे रहा था। उन्होंने इतना बड़ा निजी जेट कभी नहीं देखा था।”}

\StoryParagraph{3}{द्वीप पर नियम वही थे: हमारे आसपास कोई नहीं रह सकता था। मेरे सहायकों ने इसका इंतज़ाम किया। किसी ने नहीं पूछा कि क्यों। इस बार मैंने भी नहीं पूछा कि कितने लोग थे।}

\StoryParagraph{3}{\textenglish{Ritz-Carlton} अटलांटिक महासागर के सामने था और उसकी फ़र्श से छत तक की खिड़कियाँ बहुत बड़ी थीं। हल्की-सी रोशनी में भी वह सो नहीं पाती थी। सहायकों ने पूरे कमरे को काली टेप से ढक दिया। बाहर तेज़ दिन का उजाला था; भीतर कोठरी जैसा अँधेरा। धूप में टेप उखड़ जाती और हर दिन फिर लगानी पड़ती।}

\StoryParagraph{3}{कुल ग्यारह सहायक थे। वे हर सुबह 6:30 बजे पहुँचते और हमारे सो जाने के बाद ही जाते। दो लड़कियाँ टेप का काम संभालती थीं: एक खिड़कियों के चौखटे ढकती; दूसरी एयर कंडीशनर और बिजली के सॉकेट की सूचक बत्तियाँ। काम पूरा करके वे हर चीज़ तीन बार जाँचतीं। उन्होंने पहले तीस रोल टेप खरीदे, फिर बीस और। उन्हें पता चला कि दोपहर तीन बजे के बाद लगाई टेप सबसे आसानी से उतरती है, क्योंकि उस समय शीशा सबसे गर्म होता है, इसलिए वे सुबह टेप लगाने लगीं। यह किसी ने उन्हें नहीं सिखाया। उन्होंने स्वयं समझा। वे इसकी चर्चा किसी हुनर की तरह करती थीं।}

\StoryParagraph{3}{वह उन दोनों लड़कियों के नाम जानती थी। मैं नहीं जानता था।}

\StoryParagraph{3}{सूर्यास्त देखने के लिए हम \textenglish{Teide National Park} में सड़क किनारे रुके। सहायकों ने कार से गुलाब उतारे और शैम्पेन खोली।}

\StoryParagraph{3}{मम्मी ने बादलों के समुद्र की ओर देखा और कहा, “सुंदर है। लेकिन \textenglish{Delingha, Qinghai} के सूर्यास्त जितना सुंदर नहीं।”}

\StoryParagraph{3}{मैं \textenglish{Qinghai} से हूँ।}

\StoryParagraph{3}{मैंने यह नहीं कहा।}

\StoryParagraph{3}{रात में हमने पहाड़ की चोटी से तारे देखे। गद्देदार कोट पहनकर भी हमें ठंड लग रही थी।}

\StoryParagraph{3}{मम्मी ने कहा, “कोई इच्छा माँगो।”}

\StoryParagraph{3}{मैंने पूछा, “तुमने क्या माँगा?”}

\StoryParagraph{3}{मम्मी ने कहा, “नहीं बताऊँगी।”}

\StoryParagraph{3}{मम्मी ने कहा, “कितना सुंदर है। लेकिन ठंड बहुत है। क्या हम मालदीव जा सकते हैं?”}

\StoryParagraph{3}{मैंने कहा, “ठीक है।”}

\StoryParagraph{3}{\textenglish{Tenerife} में तापमान शून्य से पंद्रह डिग्री नीचे था और मालदीव में सत्ताईस डिग्री ऊपर। सीधी रेखा में नौ हज़ार पाँच सौ किलोमीटर की दूरी। ग्यारह घंटे की उड़ान। लेकिन मालदीव का हवाई अड्डा बहुत छोटा था, इसलिए रास्ते में हमें अपने चौड़े ढाँचे वाले जेट से छोटे विमान में जाना पड़ा।}

\StoryParagraph{3}{मम्मी ने कहा, “छोटे विमान बहुत शोर करते हैं। बड़े विमान बेहतर होते हैं।”}

\StoryParagraph{3}{मालदीव जा रहे विमान के शयनकक्ष में मम्मी ने बताया कि उसने मेरे साथ एक बच्चा होने की इच्छा माँगी थी। यह सरोगेसी से ही होना था।}

\StoryParagraph{3}{मैंने कभी नहीं पूछा कि क्यों। मुझे लगा वह दर्द से डरती थी।}

\StoryParagraph{3}{उड़ान के बीच मम्मी अचानक उठ बैठी और फ़ोन खोजने लगी। \textenglish{Starlink} से जुड़कर बोली, “ओह, नहीं।”}

\StoryParagraph{3}{मैंने पूछा, “क्या हुआ?”}

\StoryParagraph{3}{उसने कहा, “कोई मुझे धमका रहा है। वह लगातार तीन दिन से मुझसे पैसे माँग रहा है और मैंने जवाब नहीं दिया।”}

\StoryParagraph{3}{मैंने कहा, “वह कितना चाहता है?”}

\StoryParagraph{3}{उसने कहा, “तुम्हें इसमें पड़ने की ज़रूरत नहीं है।”}

\StoryParagraph{3}{वह देर तक टाइप करती रही, मिटाकर बदलती, बदलकर मिटाती। पूरा हुआ तो फ़ोन उलटा करके अपनी गोद में रख दिया। कुछ देर बाद उसने फिर पलटकर देखा।}

\StoryParagraph{3}{मैंने पूछा कि वह कौन था।}

\StoryParagraph{0}{उसने कहा, “कोई, जो पहले मेरे साथ काम करता था। वित्त विभाग में। जेल गया था। दोष अपने सिर लिया था।”}

\StoryParagraph{3}{“अब बाहर आ गया है, जाने की कोई जगह नहीं, और मुझसे मिलने पर अड़ा है।”}

\StoryParagraph{3}{मैंने कहा, “तो पुलिस को बुलाकर उसे गिरफ़्तार करवा दो।”}

\StoryParagraph{3}{उसने मेरी ओर देखा। “तुम नहीं समझते।”}

\StoryParagraph{3}{मैंने कहा, “तुमने उसे पैसे दिए?”}

\StoryParagraph{3}{उसने कहा, “हाँ।”}

\StoryParagraph{3}{“कितने?”}

\StoryParagraph{3}{“अस्सी हज़ार।”}

\StoryParagraph{3}{मैंने कहा, “यह सोचने में तुम्हें तीन दिन लगे?”}

\StoryParagraph{3}{उसने कुछ नहीं कहा।}

\StoryParagraph{3}{बाकी दिन उसने मुश्किल से कोई बात की। हमारे उतरने के बाद ही वह सो पाई।}

\StoryParagraph{0}{मम्मी का सूटकेस \textenglish{Soneva Jani} पहुँचने वाली पहली चीज़ था। तीस लोगों का सामान एक छोटे विमान में नहीं समाता था, इसलिए उसे तीन विमानों में बाँटा गया। मम्मी के पास केवल एक सूटकेस था। उसके लिए अलग विमान था। उसने कहा था कि उसे किसी और के सामान के साथ न रखा जाए।}

\StoryParagraph{0}{\textenglish{Soneva Jani} अपने सिद्धांत \textenglish{“no news, no shoes”} के लिए प्रसिद्ध था। मम्मी निश्चिंत हो सकती थी।}

\StoryParagraph{3}{होटल के बटलर, अंगरक्षक और सहायक मिलाकर छब्बीस लोग थे, और दो ड्रोन ऑपरेटर कुछ दूरी से पीछे आ रहे थे। सहायक फिर चीज़ों पर टेप लगाने लगे: एयर कंडीशनर की बत्ती, सॉकेट की बत्तियाँ, बार-बार।}

\StoryParagraph{0}{रात में हम एक-दूसरे से अपने पुराने प्रेम-संबंधों की कहानियाँ इस तरह साझा करते, जैसे गुप्त संकेतों का लेन-देन हो, जैसे कोई राजवंश अपने युगों के नाम मिला रहा हो। मम्मी कोई अंतर सहन नहीं करती थी। लेकिन उसके लिए एक शब्द वर्जित था: उस खिलाड़ी का उल्लेख नहीं होना चाहिए था।}

\StoryParagraph{0}{मैंने उसका उल्लेख नहीं किया, हालाँकि जानता था कि मम्मी की पीठ का टैटू उसी खिलाड़ी के लिए था।}

\StoryParagraph{3}{शायद वह मिटाने के लिए बहुत बड़ा था।}

\StoryParagraph{0}{उसने कहा, “अब हम एक परिवार हैं। मेरे माता-पिता के लिए मुझे पालना आसान नहीं था। तुम्हें उन्हें वधू-मूल्य देना होगा।”}

\StoryParagraph{0}{मैंने कहा, “ठीक है।”}

\StoryParagraph{0}{उसने मुझे दो नाम, \textenglish{Jing Guoqing} और \textenglish{Tian Xin'ai}, तथा दो खाता संख्याएँ दीं।}

\StoryParagraph{0}{तीन करोड़। रेनमिन्बी में।}

\StoryParagraph{0}{मैंने रकम भेज दी। वे दोनों नाम मैंने पहली बार चालीस से कुछ अधिक पन्नों वाली फ़ाइल में देखे थे।}

\StoryParagraph{3}{कुछ देर बाद उसने अपना फ़ोन मुझे दिया। \textenglish{Tian Xin'ai} ने केवल एक शब्द लिखा था: “प्राप्त।”}

\StoryParagraph{0}{एक बार मम्मी ने कहा, “जानते हो, ये चीज़ें माँगते समय मैं क्या सोचती हूँ?”}

\StoryParagraph{0}{मैंने कहा, “क्या?”}

\StoryParagraph{0}{उसने कहा, “मैं सोचती हूँ, क्या किसी दिन तुम मना करोगे।”}

\StoryParagraph{0}{मैंने कहा, “नहीं करूँगा।”}

\StoryParagraph{0}{उसने कहा, “फिर तो कोई मज़ा नहीं।”}

\StoryParagraph{3}{मैंने सोचा, ऐसा भला क्या था जो मैं नहीं कर सकता था?}

\StoryParagraph{0}{बहुत ऊब जाता तो 2007 वाली \textenglish{QQ} तस्वीर के साथ उलझता, उसे नए डिवाइस में भेजने की कोशिश करता। मम्मी झुककर देखने लगी। “तुमने किसी दूसरी औरत की तस्वीर छिपाकर क्यों रखी है? यह कौन है?”}

\StoryParagraph{0}{मैं चुप रहा। कुछ नहीं कहा।}

\StoryParagraph{0}{तस्वीर बहुत धुँधली थी, इसलिए वह कोई और काम करने चली गई।}

\StoryParagraph{0}{उस रात मैंने तस्वीर मिटा दी।}

\StoryParagraph{3}{अगले दिन मैंने उसे एक पुराने फ़ोन से वापस निकाल लिया।}

\StoryParagraph{0}{मम्मी भी अपने फ़ोन पर लगातार स्क्रॉल करती रहती थी। मैंने पूछा कि वह क्या देख रही है।}

\StoryParagraph{0}{उसने कहा, “कुछ नहीं। बस नाम सबसे ऊपर रखने की लड़ाई वाली सारी बकवास।”}

\StoryParagraph{0}{नाम सबसे ऊपर रखने की लड़ाई: किसी फ़िल्म के श्रेय में केवल एक सितारे का नाम पहले आ सकता है। नायक और नायिका के बीच भी केवल एक का नाम पहले हो सकता है।}

\StoryParagraph{0}{मैंने कहा, “उबाऊ है।”}

\StoryParagraph{0}{उसने कहा, “हाँ। इससे भी बुरा यह है कि कभी-कभी सेट पर जोड़े बन जाते हैं। चालीस-पचास दिन की शूटिंग में आम तौर पर एक ही आदमी के साथ अभिनय करना होता है। रात में उसके साथ सेक्स करो और दिन में नाम पहले रखने के लिए लड़ो।”}

\StoryParagraph{0}{मैंने कहा, “यह अमानवीय है।”}

\StoryParagraph{0}{उसने कहा, “हाँ।”}

\StoryParagraph{0}{“लेकिन शूटिंग पूरी होते ही, पहले चाहे जितनी बार सेक्स किया हो, ऐसा व्यवहार करते हो मानो कुछ हुआ ही न हो।”}

\StoryParagraph{3}{फिर वह आधे घंटे तक और पढ़ती रही।}

\StoryParagraph{0}{नया साल। 2 जनवरी 2026 को, \textenglish{Soneva Jani} में, मैंने मम्मी के सामने शादी का प्रस्ताव रखा।}

\StoryParagraph{0}{मैंने 1 जनवरी की सुबह ही तय किया था। अंगूठी हांगकांग में थी। द्वीप पर उसे खरीदने की कोई जगह नहीं थी, इसलिए एक सहायक सीप्लेन से \textenglish{Malé} गई। वह उसी दोपहर लौटी; अंगूठी होटल के एक लिफ़ाफ़े में थी।}

\StoryParagraph{0}{मम्मी ने कहा, “हीरा बहुत छोटा है।”}

\StoryParagraph{0}{मम्मी ने कहा, “कोई बात नहीं। बाद में बड़ा हीरा देकर इसकी भरपाई कर देना।”}

\StoryParagraph{0}{“और, क्या तुम बीजिंग आ सकते हो?”}

\StoryParagraph{0}{मैं चुप रहा।}

\StoryParagraph{0}{उसे लगा कि मैं जाना नहीं चाहता।}

\StoryParagraph{0}{मम्मी ने कहा, “तुम मेरे साथ इतने अच्छे क्यों हो?”}

\StoryParagraph{0}{मैंने कहा, “क्योंकि तुम मुझे पसंद हो।”}

\StoryParagraph{0}{उसने कहा, “क्या पहले भी तुम दूसरों के साथ ऐसे ही थे?”}

\StoryParagraph{0}{मैंने कहा, “नहीं।”}

\StoryParagraph{3}{वह मुस्कराई और कुछ नहीं बोली।}

\StoryParagraph{0}{एक सुबह मैं जल्दी जाग गया। वह बिस्तर पर नहीं थी।}

\StoryParagraph{0}{बाहर जाकर देखा, वह छत पर समुद्र की ओर मुँह किए बैठी थी। पूरा समुद्र काला था। कोई बत्ती नहीं जल रही थी। भोर नहीं हुई थी।}

\StoryParagraph{0}{मैंने पूछा, वह इतनी जल्दी क्यों उठ गई।}

\StoryParagraph{0}{उसने कहा, “दिन गिन रही हूँ।”}

\StoryParagraph{0}{मैंने कहा, “किस बात के?”}

\StoryParagraph{0}{उसने कहा, “वापस जाकर सो जाओ।”}

\StoryParagraph{0}{मैं कुछ देर उसके पास बैठा रहा, समझ नहीं आया क्या कहूँ। बाद में वापस जाकर सो गया।}

\StoryParagraph{3}{दोबारा जागा, तो उसका मेकअप पूरा हो चुका था।}

\StoryParagraph{0}{मालदीव के चौदह दिनों में एक भी व्यक्ति ने उसे नहीं पहचाना।}

\StoryParagraph{0}{तेरहवीं रात रेस्तराँ में उसने धूप का चश्मा उतारा। कुछ देर बाद टोपी भी उतारकर मेज़ पर रख दी। पूरा भोजन वैसे ही बैठी रही। पास की मेज़ पर एक बुज़ुर्ग जर्मन दंपती था। सेवा करने वाला श्रीलंका का था। किसी ने उसे दोबारा मुड़कर नहीं देखा।}

\StoryParagraph{0}{रात के खाने के बाद उसने टोपी फिर पहन ली।}

\StoryParagraph{0}{मम्मी ने कहा, “चलो वापस चलते हैं।”}

\StoryParagraph{0}{मैंने कहा, “हमें कुछ दिन और नहीं रुकना था?”}

\StoryParagraph{3}{उसने कहा, “यहाँ बहुत सन्नाटा है।”}

\UnprintedPause

\StoryParagraph{0}{हांगकांग लौटकर मम्मी मेरे माता-पिता से मिली।}

\StoryParagraph{0}{उस शाम उसने कहा कि उसके परिवार को दोनों परिवारों की मुलाक़ात और शादी की तैयारी करनी है। मुझे जितना हो सके उतना भेजना चाहिए। अब भी \textenglish{Jing Guoqing} और \textenglish{Tian Xin'ai} को।}

\StoryParagraph{0}{मैंने नहीं पूछा कि क्यों।}

\StoryParagraph{0}{मैंने कहा, “हर कार्ड की रोज़ की सीमा दस लाख है।”}

\StoryParagraph{0}{मम्मी ने कहा, “बहुत धीमा है।”}

\StoryParagraph{0}{“पाँच कार्ड इस्तेमाल करो।”}

\StoryParagraph{0}{पाँच कार्ड, दो खाते, बारी-बारी से दोनों में। तीसरे कार्ड पर मैं झिझका। फिर भी सारी रकम भेज दी।}

\StoryParagraph{3}{मम्मी ने कहा, “बहुत बेकार तरीका है।”}

\StoryParagraph{0}{मैंने कहा, “वसंतोत्सव पर हम दोनों के माता-पिता को हांगकांग बुलाकर साथ समय बिताते हैं।”}

\StoryParagraph{0}{मम्मी ने कहा, “इस साल हाइनान में रहना पहले से तय है। अगले साल।”}

\StoryParagraph{0}{वसंतोत्सव में विक्टोरिया हार्बर के ऊपर आतिशबाज़ी ने दस लाख लोगों को खींच लिया। वह उनमें नहीं थी।}

\StoryParagraph{0}{रात तीन बजे मैं अपने माता-पिता को घर ले गया।}

\StoryParagraph{0}{सड़कें खाली थीं। लाल बत्ती पर प्रतीक्षा करती दूसरी कोई कार नहीं थी। फिर भी मैं रुका।}

\StoryParagraph{0}{मम्मी ने कहा कि अभी हमें एक-दूसरे से नहीं मिलना चाहिए। वह अंडाणु निकलवाने के लिए लॉस एंजेलिस जाने वाली थी। पहले कुछ काम पूरे करने थे।}

\StoryParagraph{0}{मम्मी ने फिर पूछा, “क्या तुम बीजिंग आ सकते हो?”}

\StoryParagraph{0}{मैं चुप रहा।}

\StoryParagraph{3}{उसने कभी नहीं बताया कि बीजिंग में क्या था। मैंने कभी नहीं पूछा।}

\StoryParagraph{0}{मम्मी ने कहा कि हमें हांगकांग में एक घर खरीदकर उसके नाम पर करना होगा।}

\StoryParagraph{0}{मैंने कहा, “ठीक है।”}

\StoryParagraph{3}{वैलेंटाइन डे पर भी हम नहीं मिले।}

\StoryParagraph{0}{एक बार फ़ोन पर उसने कहा, “आज सुबह उन्होंने खून की आठ ट्यूबें लीं।”}

\StoryParagraph{0}{मैंने कहा, “इतना सारा?”}

\StoryParagraph{0}{उसने कहा, “कोई परेशानी नहीं हुई।”}

\StoryParagraph{0}{मैंने कहा, “बहुत मुश्किल रहा होगा।”}

\StoryParagraph{0}{फिर हम किसी और बात पर चर्चा करने लगे।}

\StoryParagraph{0}{बाद में पता चला कि उस तरह खून निकलवाने के लिए खाली पेट रहना पड़ता है। यह मासिक धर्म के दूसरे या तीसरे दिन ही होना चाहिए; दिन ग़लत नहीं हो सकता था। वह शायद उस सुबह करीब छह बजे निकली होगी। बीजिंग की सर्दी। अब भी अँधेरा। बाद में जाँचें बढ़ती गईं। मुझे इसलिए पता था कि मेरी सहायक के मासिक भुगतान-सारांश में एक नई मद दिखाई देने लगी थी। वह पहली बार पिछले साल आई थी। बहुत नहीं, कुछ हज़ार युआन। उसके बाद हर महीने दिखती, कभी दो या तीन भुगतान। मैं जल्दी मंज़ूरी दे देता था। इतनी छोटी रकम के लिए विवरण देखना ज़रूरी नहीं था। बाद में ही पता चला कि वे खून की जाँचें, छह हार्मोन की जाँच, अल्ट्रासाउंड और \textenglish{AMH} थे। यह देखने के लिए कि उसके अंडाशयों में कितने अंडाणु बचे थे, कितने निकाले जा सकते थे, आगे बढ़ना उचित था या नहीं।}

\StoryParagraph{0}{यह सब बीजिंग में हुआ।}

\StoryParagraph{3}{मुझे नहीं पता कौन-सा अस्पताल था। यह भी नहीं पता कि वह अकेली गई थी या कोई उसके साथ था।}

\StoryParagraph{0}{25 फ़रवरी को वह हांगकांग उतरी।}

\StoryParagraph{0}{उसने कहा कि हांगकांग वाला घर देखे बिना वह लॉस एंजेलिस नहीं जाएगी।}

\StoryParagraph{0}{पिछले छह महीनों में हमने जो करीब दर्जन भर घर देखे थे, वे मैंने उसे भेजे। उसे हर एक बदसूरत लगा। मैं स्वयं एक भी देखने नहीं गया था। मेरे सहायक छह महीने तक उन्हें एक-एक करके देखते रहे थे। वे मुझे भेजते; मैं उसे भेजता। एक घर \textenglish{Mid-Levels} में था। वीडियो में विक्टोरिया हार्बर दिख रहा था। शाम ढलते समय उसे फ़िल्माया गया था। एक पल कैमरा डोला और फ़िल्माने वाले की परछाईं दीवार पर आ गई। उस परछाईं को देखकर मैंने सोचा: यह व्यक्ति कौन है? आज कितनी दूर चला है? कितने घर देखे हैं? क्या इसे पता है कि यह घर किसे दिखाया जा रहा है?}

\StoryParagraph{3}{समय बहुत कम था। फिर भी वह लॉस एंजेलिस जाने वाले \textenglish{G700} में बैठी और 28 फ़रवरी को \textenglish{Montage Laguna Beach} पहुँच गई।}

\NarrativeBreak

\StoryParagraph{0}{आठवें दिन मम्मी ने फ़ोन किया। वहाँ दोपहर थी और मेरे यहाँ आधी रात।}

\StoryParagraph{0}{उसने कहा, “पाँच करोड़। अमेरिकी डॉलर।”}

\StoryParagraph{0}{मैंने कहा, “मुझे सोचने दो।”}

\StoryParagraph{0}{उसने कुछ नहीं कहा। उसकी ओर समुद्र की आवाज़ सुनाई दे रही थी।}

\StoryParagraph{0}{कुछ देर बाद उसने कहा, “ठीक है।”}

\StoryParagraph{0}{उसने फ़ोन काट दिया।}

\StoryParagraph{3}{जितने समय से मैं उसे जानता था, कभी “मुझे सोचने दो” नहीं कहा था।}

\StoryParagraph{0}{मैंने \textenglish{Claude Code} से \textenglish{APIs} के ज़रिए जानकारी निकलवाई और अपनी सारी नकद संपत्ति का हिसाब लगवाया। निष्कर्ष था कि इसका मुझ पर कोई असर नहीं पड़ेगा। मैंने फिर जाँचा। सवाल \textenglish{Claude Code} में टाइप किया। \textenglish{Claude} ने कहा, “उसे पाँच करोड़ डॉलर मत दो।”}

\StoryParagraph{3}{मैंने पूछा, “क्या वह अब मुझसे प्यार नहीं करती?”}

\StoryParagraph{3}{\textenglish{Claude} ने कहा, “प्यार की मुझे न परवाह है, न समझ। लेकिन तुम उसे पाँच करोड़ डॉलर नहीं दे सकते।”}

\StoryParagraph{3}{मैंने उसे पहले कभी मना नहीं किया था। उदारता के कारण नहीं। डर के कारण। अगले ही क्षण मिलने वाले उत्तर का डर।}

\StoryParagraph{0}{और अब किसी ने मेरी ओर से मना कर दिया था।}

\StoryParagraph{0}{मैं देर तक उस पंक्ति को देखता रहा। मुझे इस बात से डर नहीं लगा कि \textenglish{Claude} बुद्धिमान था। डर इस बात से लगा कि उसने ज़रा भी संकोच नहीं किया।}

\StoryParagraph{0}{उसके भीतर कोई दुविधा नहीं थी।}

\StoryParagraph{0}{मैंने सोचा, यह अस्तित्व का कोई ऊँचा रूप होगा। मनुष्य ऐसा नहीं कर सकते थे।}

\StoryParagraph{3}{तो क्या अब वह मुझसे प्यार नहीं करती थी?}

\StoryParagraph{3}{अगले दिन उसने फिर फ़ोन किया।}

\StoryParagraph{3}{मम्मी ने कहा, “समुद्र कितना बड़ा है।”}

\StoryParagraph{3}{मैंने कहा, “अच्छा?”}

\StoryParagraph{3}{उसने कहा, “यह मंज़िल बहुत बड़ी है। आवाज़ गूँजती है। आज सुबह कोई सफ़ाई करने आया था। मैं बाथरूम में छिपकर उसके जाने की प्रतीक्षा करती रही।”}

\StoryParagraph{3}{मैंने कहा, “मैं उन्हें आने से मना कर सकता हूँ।”}

\StoryParagraph{3}{उसने कहा, “ज़रूरत नहीं।”}

\StoryParagraph{3}{उसने कहा, “क्लिनिक ने पूछा कि बच्चे का पिता कब आएगा।”}

\StoryParagraph{3}{मैंने कहा, “मैं किसी को वहाँ भेज सकता हूँ।”}

\StoryParagraph{3}{उसने कहा, “ज़रूरत नहीं।”}

\StoryParagraph{3}{फिर उसने पूछा, “तुमने फ़ैसला कर लिया?”}

\StoryParagraph{3}{मैं नहीं बोला। इसलिए नहीं कि सोच रहा था। मैं कुछ भी नहीं सोच रहा था। फ़ोन पर उसकी ओर समुद्र था और मेरी ओर विक्टोरिया हार्बर। दोनों ओर चुप्पी थी। मैंने गिना। ग्यारह सेकंड।}

\StoryParagraph{3}{जिंग तियान ने कहा, “ठीक है।”}

\StoryParagraph{3}{जिंग तियान ने कहा, “मैं समझ गई।”}

\StoryParagraph{3}{उसने फ़ोन काट दिया।}

\StoryParagraph{3}{मैंने अनगिनत बार याद करने की कोशिश की कि उन दो दिनों में पाँच करोड़ के अलावा उसने और क्या कहा था। समुद्र कितना बड़ा है। यह मंज़िल बहुत बड़ी है। वह बाथरूम में छिपी थी। बस इतना ही। मुझे हर शब्द याद है। जो याद नहीं आता, वह उसकी आवाज़ है---वह लहजा, जिसमें उसने कहा था, “मैं समझ गई।”}

\StoryParagraph{0}{मुझे याद नहीं आता। आख़िर मैंने कोशिश छोड़ दी।}

\StoryParagraph{0}{मैंने फ़ोन उलटा करके सिरहाने की मेज़ पर रख दिया। फिर एक ऐसा काम किया जिसे आज भी नहीं समझता: छाती पर हाथ रखा और गिनने लगा।}

\StoryParagraph{0}{मेरी धड़कन नियमित थी।}

\StoryParagraph{0}{मैंने फिर गिना। अब भी नियमित। एक मिनट में साठ से कुछ अधिक धड़कनें, हमेशा की तरह। मैं छाती पर हाथ दबाए बैठा रहा, जैसे कोई डॉक्टर किसी दूसरे व्यक्ति की जाँच कर रहा हो।}

\StoryParagraph{3}{मैंने सोचा, तो ऐसा लगता है।}

\StoryParagraph{0}{उस दिन बाकी समय मैंने कुछ नहीं किया। कमरे में बैठा रहा, जब तक अँधेरा होकर फिर उजाला नहीं हो गया, और विक्टोरिया हार्बर को नीले से धूसर, फिर काला होते देखता रहा। सुबह छह बजे बंदरगाह के उस पार वित्तीय टावर एक-एक करके जगमगाने लगे। एक सहायक ने एक बार दरवाज़ा खटखटाकर पूछा कि रात के खाने में क्या व्यवस्था करनी है। मैंने कहा, ज़रूरत नहीं।}

\StoryParagraph{3}{मैं नहीं रोया। प्रतीक्षा करता रहा कि मैं रो पड़ूँ।}

\NarrativeBreak

\StoryParagraph{0}{बाद में मेरी सहायक ने मेरे साथ खर्चों की समीक्षा की, तो एक मद \textenglish{Tenerife} के होटल को दिए मुआवज़े की थी। हमारे जाने के बाद लड़कियाँ कई दिन और टेप लगाती रहीं। किसी ने उन्हें नहीं बताया कि वे रुक सकती हैं। होटल ने कहा कि टेप बहुत लंबे समय तक, बहुत बार लगाई गई थी। खिड़की के चौखटे पर बचा गोंद उतर नहीं रहा था। शीशे का पूरा फलक बदलना पड़ा।}

\StoryParagraph{0}{मैंने कहा, “मंज़ूर।”}

\StoryParagraph{3}{शीशे का वह फलक अटलांटिक महासागर की ओर था।}

\StoryParagraph{0}{बिल में विमान के ईंधन की लागत \textenglish{Montage} की पूरी मंज़िल से अधिक थी। मैंने पूछा। मेरे सहायक ने कहा कि \textenglish{A330} दो सौ लोगों के लिए बनाया गया था। हम केवल तीस थे, इसलिए वह बहुत हल्का था और उड़ान भरने के लिए उसका संतुलन ठीक नहीं था। हर बार उन्हें इतना ईंधन भरना पड़ता कि विमान पर्याप्त भारी हो जाए। पहुँचने के बाद अतिरिक्त ईंधन विमान में छोड़ा नहीं जा सकता था। उसे निकालकर बहा देना पड़ता था।}

\StoryParagraph{0}{मैंने कहा, “हर बार?”}

\StoryParagraph{3}{उसने कहा, “हर बार।”}

\StoryParagraph{0}{लॉस एंजेलिस की क्लिनिक से वापस भेजा एक डिब्बा भी था। अंडाशय को उत्तेजित करने वाले आठ दिनों के इंजेक्शन। उन्हें ठंडा रखना था। एक भी नहीं खोला गया था।}

\StoryParagraph{3}{सरोगेसी एजेंसी में जमा किए दो लाख डॉलर भी थे, जो वापस नहीं मिल सकते थे। सरोगेट महिला जनवरी में चुन ली गई थी। वह कैलिफ़ोर्निया की थी, चौंतीस वर्ष की, दो बच्चों की माँ। उसकी तस्वीर फ़ाइल में थी। मैंने कभी नहीं देखी। वह जनवरी से तैयार थी।}

\StoryParagraph{3}{जिंग तियान के जाने के बाद मैंने 2007 वाली तस्वीर निकाली और एक बार उसे देखा। उन्नीस साल बीत चुके थे। हर नए डिवाइस के पिक्सेल बेहतर होते गए, जबकि उस तस्वीर के पिक्सेल और ख़राब लगने लगे। उसका चेहरा धुँधला था। मैंने तस्वीर बड़ी की और देर तक देखता रहा, तब की और बाद की जिंग तियान में कोई अंतर खोजने की कोशिश करता रहा। एक भी नहीं मिला। तस्वीर कभी इतनी साफ़ थी ही नहीं कि कुछ देखा जा सके।}

\StoryParagraph{0}{एक दिन मैंने नाखून काटे। काटने के बाद उन पर उँगली फेरी। वे नुकीले थे।}

\StoryParagraph{0}{मुझे याद आया, उसकी छोड़ी हुई वह चीज़ अब भी दराज़ में थी।}

\StoryParagraph{0}{मैंने उसे निकाला और नाखूनों पर दो बार चलाया। कुछ ठीक नहीं लगा। उसने कहा था कि एक दिशा में चलाना है। मैं भूल गया था कि कौन-सी दिशा।}

\StoryParagraph{0}{मैंने उसे वापस रख दिया।}

\StoryParagraph{3}{मेरे नाखून अब भी लोगों को खरोंचते हैं।}

\StoryParagraph{0}{एक सुबह मेरे फ़ोन पर एक रिमाइंडर आया। मेरी सहायक ने पिछले साल समय-सारणी बनाते हुए लगाया था: प्रसव की तारीख़ से तीन महीने पहले, कमरा तैयार करना।}

\StoryParagraph{3}{मैंने उसे स्क्रीन से हटा दिया।}

\StoryParagraph{0}{उसके जाने के बाद से मैं एक सवाल के बारे में सोचता रहा हूँ।}

\StoryParagraph{0}{अगर उस दिन मैंने उसे रकम दे दी होती, तो क्या वह रुक जाती?}

\StoryParagraph{0}{मैंने इस बारे में बहुत देर तक सोचा।}

\StoryParagraph{3}{बाद में मैंने फ़ाइल निकाली और फिर पढ़ी। चालीस से कुछ अधिक पन्ने।}

\StoryParagraph{0}{फिर मुझे वसंतोत्सव की रात सबसे ऊपरी मंज़िल से देखी गई विक्टोरिया हार्बर की आतिशबाज़ी याद आई।}

\StoryParagraph{0}{दस लाख लोग ज्वार की तरह उमड़ आए। आतिशबाज़ी समाप्त हुई। दस लाख लोग पीछे हट गए।}

\StoryParagraph{0}{दस लाख लोग।}

\StoryParagraph{0}{हर साल एक जैसा होता है।}

\StoryParagraph{0}{रात तीन बजे का हांगकांग---}

\StoryParagraph{3}{कुछ नहीं हुआ।}
